\chapter{SWHID tool implementations}
\label{app:swhid-tools}

A \textsc{swhid} is, by design, something anyone can compute: the whole point of an intrinsic
identifier is that it owes nothing to a registry, so a name you can only obtain from a central
service would betray the idea. Now that the \textsc{swhid} is an international standard
(ISO/IEC~18670:2025), a small ecosystem of independent implementations has grown up around it ---
in several languages, at several levels of ambition --- and they all compute, for the same bytes,
the same name. That is the standard doing its job.

For everyday work you need only two of them. The \emph{reference implementation} is
\texttt{swhid-rs}, written in Rust and released under the \textsc{mit} licence; it tracks the ISO
specification directly and is the tool this note reaches for first. The Python
\texttt{swh.model} package (the \texttt{swh identify} command) remains the most complete in
coverage --- snapshots, deposits, the full archive toolchain --- and \texttt{swh scanner} answers
the complementary question of \emph{what the archive already holds}. The rest of the table records
the alternatives you may meet in the wild, and the no-install resolvers that turn a \textsc{swhid}
back into the code it names.

\begingroup
\small
\setlength{\extrarowheight}{2pt}
\begin{table}[htbp]
  \centering
  \caption[SWHID tool implementations]{Implementations and resolvers for the \textsc{swhid}
  standard (ISO/IEC~18670). The Rust tool \texttt{swhid-rs} is the reference implementation; the
  Python \texttt{swh.model} CLI is the most feature-complete. Resolvers need no installation at
  all.}
  \label{tab:swhid-tools}
  \begin{tabularx}{\linewidth}{@{}>{\raggedright\arraybackslash}p{2.5cm}%
                                  >{\raggedright\arraybackslash}p{2.0cm}%
                                  >{\raggedright\arraybackslash}X@{}}
    \toprule
    \textbf{Name} & \textbf{Language /\newline ecosystem} & \textbf{What it does \& how to use it} \\
    \midrule
    \texttt{swhid-rs}
      & Rust \newline (\textsc{mit})
      & \emph{Reference implementation} of ISO/IEC~18670. Computes and verifies core
        \textsc{swhid}{s}. \texttt{cargo install swhid} (add \texttt{-{}-features git} for
        revision support), or grab a pre-built binary from the releases page.
        \texttt{swhid content -{}-file F}, \texttt{swhid dir DIR}, \texttt{swhid verify PATH SWHID},
        \texttt{swhid parse "swh:1:\textellipsis"}.
        \scriptsize\url{https://github.com/swhid/swhid-rs} \\
    \addlinespace
    \texttt{swh.model} \newline (\texttt{swh identify})
      & Python
      & The most complete CLI: content, directory, revision \emph{and} snapshot identifiers, plus
        \texttt{-{}-verify}. \texttt{pip install 'swh.model[cli]'}, then e.g.\
        \texttt{swh identify -{}-type snapshot repo.git/}.
        \scriptsize\url{https://gitlab.softwareheritage.org/swh/devel/swh-model} \\
    \addlinespace
    \texttt{swh scanner}
      & Python
      & Walks a working tree and tells you which files and directories are \emph{already archived}
        in Software Heritage --- the complement to computing a name.
        \texttt{pip install swh.scanner}; \texttt{swh scanner scan ./repo}.
        \scriptsize\url{https://gitlab.softwareheritage.org/swh/devel/swh-scanner} \\
    \addlinespace
    \emph{updateswh} \newline (Permalinks)
      & Browser \newline extension
      & The point-and-click route: a button on GitHub / GitLab / Bitbucket pages that triggers
        archival and then exposes the artifact's \textsc{swhid} via the archive's \emph{Permalinks}
        tab. No local install of a hashing tool needed.
        \scriptsize\url{https://www.softwareheritage.org/browser-extensions/} \\
    \addlinespace
    Archive \& \texttt{n2t.org} resolvers
      & Web service \newline (no install)
      & Turn a \textsc{swhid} \emph{back} into code: paste it after
        \texttt{https://archive.softwareheritage.org/} to browse the bytes, or resolve via the
        \texttt{n2t.org} name-resolver (\texttt{https://n2t.net/}).
        \scriptsize\url{https://archive.softwareheritage.org/} \\
    \addlinespace
    \texttt{swhid} for OCaml
      & OCaml
      & A community library (OCamlPro) to parse, build and manipulate \textsc{swhid}{s}.
        \texttt{opam install swhid}.
        \scriptsize\url{https://github.com/OCamlPro/swhid} \\
    \addlinespace
    \texttt{swhid-go}, \texttt{swhid-ruby}
      & Go, Ruby
      & Independent ports of the standard for Go and Ruby toolchains.
        \scriptsize\url{https://github.com/andrew/swhid-go} \\
    \bottomrule
  \end{tabularx}
\end{table}
\endgroup

\noindent A practical rule of thumb: reach for \texttt{swhid-rs} when you want the
reference behaviour and a single fast binary; fall back to \texttt{swh identify} for snapshots and
the deeper archive workflow; and remember that for a one-off check you need install nothing at all
--- the resolver at \url{https://archive.softwareheritage.org/} will take any \textsc{swhid} and
hand you back the code it names.
